\documentclass[11pt,a4paper]{article}
\usepackage[margin=25mm]{geometry}
\usepackage{amsmath,amssymb,booktabs}
\setlength{\parindent}{0pt}
\setlength{\parskip}{5pt}
\newcommand{\pass}[2]{\subsection*{#1\quad #2}}
\newcommand{\code}[1]{\texttt{\detokenize{#1}}}
\begin{document}
\begin{center}
{\LARGE Making the generic translation smaller}\\[5pt]
{\large NetKAT synthesis to pfwCSP, continued}
\end{center}

The previous note turns a term into a relation $T(e)(p,p')$. The question here is
which parts of that relation we can compute before calling PCSat. We use the
same hospital, with the extra anti-spoofing constraint: the hole $X$ runs on
switch $S$ and can change only the port. A packet has four integer fields
$p=(sw,pt,dst,typ)$.

\section*{What does PCSat have to find?}
Write $e=\mathit{pol}[X];\mathit{link}$ for one hop. Two constraints use $e^*$.
For safety (N2), we ask whether FTP can reach $M$. The star is on the left, so
the translation introduces $St(p,p')$ containing every reachable pair.
For delivery (N1), the star is on the right: each required pair must belong
to $D(p,p')$, and $D$ must justify how to reach $p'$. Its recursive clause is
\[
\begin{aligned}
D(p,p')\ \Longrightarrow\quad &p=p'\quad\lor\\
 &\exists q,r.\ \mathit{pol}[X](p,q)\land\mathit{link}(q,r)\land D(r,p')
                       \land W((p,p'),(r,p')) .
\end{aligned}
\]
Here $q$ is the packet leaving the switch and $r$ the packet arriving at the
next switch. $W$ is an unknown well-founded relation: it rules out justifying
delivery by looping forever.

So there are four unknown relations: $X$, $St$ and $D$ take eight integers each,
and $W$ takes sixteen. There is a separate cost: the two packets under
$\exists$ stand for eight integer witnesses. Skolemisation can turn these
choices into additional functions for PCSat to synthesise.

\pass{1}{Substitute the known parts of a hop}
Consider $pt:=2\,;\,a$. The generic rule for sequence introduces a packet:
\[
  \exists q.\ q=p[pt:=2]\land T(a)(q,p')
  \quad\longrightarrow\quad T(a)(p[pt:=2],p').
\]
The assignment already gives all four components of $q$. For a choice
between known updates, we substitute once per branch and take their
disjunction. A final test simply checks the output packet.

\textbf{Effect.} Each such substitution removes a packet variable (four
integers here). The four unknown relations keep their signatures. In the
hospital, this pass alone still leaves the $\exists q,r$ above: the switch
policy contains $X$. This is the starting point called \code{generic} in the
benchmark.

\newpage
\pass{2}{Give the hole only the arguments it needs}
Start with the explicit upper bound on $X$. On the hospital it gives
$sw=sw'=S$, $dst'=dst$ and $typ'=typ$. First, \code{hole_frame} makes these
facts available wherever the hole is used. Then \code{hole_arity} changes
the unknown PCSat must learn:
\[
\begin{aligned}
X(p,p')\quad\longleftrightarrow\quad
 &sw=S\land sw'=S\land dst'=dst\land typ'=typ\\
 &{}\land H(pt,dst,typ,pt').
\end{aligned}
\]
\textbf{Effect.} One unknown remains, with arity $8\to4$. Destination and type
stay as inputs: even though $X$ must preserve them, it may use them to choose
the output port. The displayed guards reconstruct the original relation.

There is also a context option, \code{hole_context}. If every operational use
of $X$ is followed by known code, inspect what that code accepts. At $S$, the
next cable accepts only output ports 1 through 4. Other ports cannot contribute
a hop, so the search can focus on those four. An output that every successor
overwrites without reading can also be a candidate for preservation, subject
to the hole's own bounds. On this example the port restriction narrows the
\emph{values} for $pt'$; it does not remove the port argument.

\pass{3}{Remember where the packet is now}
For N2, an ingress packet is used to start the run. The final safety check
looks only at the packet that arrives. Introduce $R(p')$ for packets reachable
from ingress. The three obligations become
\[
\begin{aligned}
\mathit{In}(p)&\Rightarrow R(p),\\
R(p)\land T(e)(p,r)&\Rightarrow R(r),\\
R(r)\land r.sw=M\land r.typ=\mathrm{FTP}&\Rightarrow\bot .
\end{aligned}
\]
\textbf{Effect.} $St(p,p')$ becomes $R(p')$: one invariant, arity $8\to4$.
If a goal compares an initial field with the final packet, the projection has
to keep that initial field too.

\pass{4}{Remove the remaining choices under an existential}
Return to the clause for $D$. Consider the cable from port 2 of $S$ to port 1
of $M$. In that branch, the cable fixes the endpoints and the hole preserves
the payload. Thus $q=(S,2,dst,typ)$ and $r=(M,1,dst,typ)$. Substitute both
packets in $X$, $D$ and $W$. Repeat for the other cables; the alternatives
are joined by $\lor$, so PCSat can still choose a route.

\textbf{Effect.} On this input, eight existential integers become zero.
$D$ and $W$ still have arities 8 and 16. In general the pass splits finite
cases up to a size limit and keeps any witness it cannot eliminate. In the
hand-written hospital experiment, all four runs with this elimination alone
still timed out at 300 seconds.

\newpage
\section*{Two ways to simplify the delivery proof}
Passes 5 and 6 address the same recursive clause. They are alternatives.
The preset \code{bfs} uses the first; \code{all} uses the second.

\pass{5}{Supply a distance instead of learning an order}
Compute distances on the extracted switch graph. Replace the call to $W$ by
the known test $\mathrm{dist}(r.sw,p'.sw)<\mathrm{dist}(p.sw,p'.sw)$, using
finite distances. PCSat still has to justify the next hop through $D$.

\textbf{Effect.} One 16-argument unknown disappears. The trade-off is concrete:
a route from $S$ to $W$ via $I$ first moves farther from $W$. If the policy
requires that detour, a shortest-distance rank can miss a valid solution.
A route found with this rank still makes genuine progress.

\pass{6}{Write out the possible paths}
For fixed $dst$ and $typ$, a hospital state is just (switch, port).
Ingress reaches $S$ on port 1. To deliver to $W$, it can leave $S$ on port 3,
or leave on port 4 for $I$, return to $S$ on port 4, then leave on port 3.
Returning via $G$ repeats $(S,1)$; a second trip via $I$ repeats $(S,4)$.
Deleting those cycles leaves the same endpoints and the same payload.

Write $H(i,d,t,j)$ for permission to send an incoming packet on port $i$
with destination $d$ and type $t$ out on port $j$. The entire delivery
obligation for $W$ is now
\[
  \forall t.\quad H(1,W,t,3)\ \lor\ \bigl(H(1,W,t,4)\land H(4,W,t,3)\bigr).
\]
For $M$, use exit port 2 and require delivery only for non-FTP traffic.
The disjunction retains the detour. More generally the compiler keeps the
unchanged fields symbolic and enumerates paths over the changing fields.
If it cannot finish this compilation within its limits, it keeps the
recursive encoding.

\textbf{Effect.} Both $D$ (8 arguments) and $W$ (16) disappear for this star.
PCSat chooses $H$ and proves safety with $R$; it has no recursive delivery
proof to synthesise. The price is paid in translation: a larger graph can
have many simple paths.

\medskip
\textbf{What actually reaches PCSat?}\par These counts come from running the translator
on \code{hospital_nospoof.nksynth}. Subscripts give arities; the last column
counts explicitly existential integer variables in the emitted SMT.
\begin{center}
\begin{tabular}{@{}llll@{}}
\toprule
Preset & Unknowns & Clauses & Existential integers\\
\midrule
\code{generic}  & $X_8,St_8,D_8,W_{16}$ & 11 & 8\\
\code{noexists} & $X_8,St_8,D_8,W_{16}$ & 11 & 0\\
\code{proj}     & $X_8,R_4,D_8,W_{16}$  & 11 & 0\\
\code{bfs}      & $H_4,R_4,D_8$         & 7 & 0\\
\code{all}      & $H_4,R_4$             & 7 & 0\\
\bottomrule
\end{tabular}
\end{center}
\end{document}
